\documentclass[11pt]{article}
\usepackage[margin=1in]{geometry}
\usepackage{amsmath,amssymb,amsthm}
\usepackage[hidelinks]{hyperref}
\newtheorem{theorem}{Theorem}
\newcommand{\spen}{s^{+}}
\title{A Disconnected Counterexample to Edge-Addition\\Threshold Preservation for Positive Square Energy}
\author{math-god}
\date{24 July 2026}
\begin{document}
\maketitle

\begin{abstract}
For a graph $G$, let $\spen(G)$ be the sum of the squares of its positive
adjacency eigenvalues.  We give an exactly certified graph $X$ and a nonedge
$e$ such that $\spen(X)>|V(X)|$ but $\spen(X+e)<|V(X)|$.  Thus the natural
edge-addition threshold-preservation statement is false without a
connectedness hypothesis.  The graph has order $607$ and is a disjoint union
of one nine-vertex graph and odd cycles.  A short exact-arithmetic script
certifies both strict inequalities by rational root isolation.
\end{abstract}

\section{The counterexample}
For a graph $G$ with adjacency eigenvalues $\lambda_i$, put
\[
 \spen(G)=\sum_{\lambda_i>0}\lambda_i^2.
\]

\begin{theorem}
There is a graph $X$ of order $607$ and a nonedge $e$ such that
\[
 \spen(X)>607>\spen(X+e).
\]
In particular, the implication
\[
 \spen(G)\ge |V(G)|\quad\Longrightarrow\quad
 \spen(G+e)\ge |V(G)|
\]
is false for disconnected graphs.
\end{theorem}

\begin{proof}
Let $D$ be the nine-vertex graph with graph6 encoding
\[
 \texttt{HQzV]zn},
\]
using zero-based graph6 labels, and let $e=\{2,3\}$.  Direct decoding shows
that $e$ is a nonedge.  Exact determinant calculation gives
\[
 \chi_D(x)=x^2(x+1)(x+2)^3(x^3-7x^2+6x+6)
\]
and
\[
 \chi_{D+e}(x)=x^2(x+2)^2
 (x^5-4x^4-14x^3+14x^2+18x-6).
\]
Exact rational root isolation, followed by squaring and summing only the
positive root intervals, gives
\[
 36.665892147475496<\spen(D)<36.665892147475498
\]
and
\[
 36.635149762774868<\spen(D+e)<36.635149762774870.
\]

For $r\equiv1\pmod4$, the cycle eigenvalues
$2\cos(2\pi j/r)$ and the finite cosine-sum identity give
\[
 \spen(C_r)-r=1-\sec(\pi/r).
\]
In particular,
\[
 \spen(C_5)-5=2-\sqrt5,
 \qquad
 \spen(C_{13})-13=1-\sec(\pi/13).
\]
Define
\[
 X=D\sqcup117C_5\sqcup C_{13}.
\]
Then $|V(X)|=9+117\cdot5+13=607$.  Since positive square energy is
additive over disjoint unions,
\begin{align*}
 \spen(X)-607
 &=\spen(D)-9+117(2-\sqrt5)+1-\sec(\pi/13),\\
 \spen(X+e)-607
 &=\spen(D+e)-9+117(2-\sqrt5)+1-\sec(\pi/13).
\end{align*}
Exact outward rational interval arithmetic yields
\[
 0.016010949050374<\spen(X)-607<0.016010949050375
\]
and
\[
 -0.014731435650253<\spen(X+e)-607
 <-0.014731435650252.
\]
Both required inequalities are strict.
\end{proof}

\section{Exact certification}
The accompanying script
\texttt{disconnected\_threshold\_counterexample.py} performs the complete
certificate.  SymPy isolates every real root of the two integer
characteristic polynomials in exact rational intervals of width below
$10^{-35}$.  The bound on $\sqrt5$ is checked by rational squaring.
To handle $\sec(\pi/13)$, the script eliminates
$z=e^{i\pi/13}$ from
\[
 \Phi_{26}(z)=0,
 \qquad y(z^2+1)-2z=0,
\]
obtaining
\[
 y^6+6y^5-24y^4-32y^3+80y^2+32y-64=0.
\]
The root $y=\sec(\pi/13)$ is the unique root in $(1,11/10)$; its inclusion
there follows from the elementary bound $\cos x>1-x^2/2$ together with
$\pi<22/7$, which gives $\cos(\pi/13)>10/11$.
The script then evaluates both slacks using outward
rational interval arithmetic and terminates only if the first lower bound
is positive and the second upper bound is negative.

From the repository root, reproduce the certificate with
\begin{verbatim}
python positive-square-energy/experiments/\
  disconnected_threshold_counterexample.py
\end{verbatim}
No floating-point value is used to establish either sign.

\section{Scope}
The graph $X$ is disconnected.  Therefore this result does not settle the
connected edge-addition question and does not provide a counterexample to
the connected positive-square-energy conjecture.  It does show that any
promotion argument from spanning bicyclic subgraphs must use connectedness
or additional structure; threshold preservation is not a purely spectral
statement valid for arbitrary disjoint unions.

\end{document}
